\documentclass[11pt,a4paper]{article}
\usepackage[margin=24mm]{geometry}
\usepackage{amsmath,amssymb,amsthm,hyperref}
\newtheorem{theorem}{Theorem}\newtheorem{lemma}[theorem]{Lemma}\newtheorem{corollary}[theorem]{Corollary}
\newcommand{\Sz}{\operatorname{Sz}}\newcommand{\Aut}{\operatorname{Aut}}\newcommand{\gen}[1]{\langle#1\rangle}
\title{Generating Suzuki triples by excluding proper subfields}
\author{Human Compression project: L015 working proof}\date{}
\begin{document}\maketitle
\noindent This is a private, unreviewed uniform repair of the generation step
in Leemans--Stokes--Tranchida, arXiv:2311.13522v1 (22 November 2023),
physical pages 6--8. It is not a complete L015 recovery or an assertion that
its unrestricted Theorem 3.2 is true. The accompanying complete application
supplies the source's incidence conclusions and records its corrected routes.
No finite search, new conjecture answer, formalization or factor is used here.

\section{Explicit retained structure}
Let $G=\Sz(2^n)$, $n>1$ odd, in its standard four-dimensional matrix
model and ovoid action on $2^{2n}+1$ points. Write $F$ for coefficient
squaring and $\theta=2^{(n+1)/2}$. We retain the classical Suzuki structure
used in the accepted Topic D treatment: the doubly transitive ovoid action;
unique fixed point of every involution; point Sylow $2$-group $U$, with
involutions its elementary abelian center $Z$; $C_G(z)=U$ for $1\ne z\in Z$;
and the split two-point torus acting transitively and regularly on $Z\setminus1$.
The maximal subgroups are Borels, odd cyclic torus normalizers with cyclic
Sylow $2$-subgroups, and conjugates of the canonical subfield groups
$S_d=\Sz(2^d)$, $d\mid n$, $d<n$. The first family, omitted from printed
Theorem 2.1, is included here. Nested standard subfields have nested
matrix and ovoid models; their twists agree because every extension degree
is odd. The coefficient-fixed group for $F^a$ is the corresponding
standard subfield group (also $S_1=\Sz(2)$ of order $20$).

The nonidentity odd-order elements belong to cyclic tori of orders
$q-1,q+\sqrt{2q}+1,q-\sqrt{2q}+1$, $q=2^n$. Their centralizers are those
tori, and the normalizer quotients have orders $2,4,4$, acting freely on
nonidentity torus elements. Every nonidentity element of even order belongs
to a conjugate of $U$ and has order two or four. Its eigenvalues in
characteristic two are all one, so its four-dimensional trace is zero.
These are retained classical structure inputs, not new
classification results. These exact retained interfaces, with their source
locators and earlier Topic D provenance, are listed in the accompanying
dependency record. No proof in another application is silently imported:
the local trace, subfield-normalizer and counting arguments are written here.
The complete L122 common paper remains a separate, fully charged dependency
of the incidence application. Complete external structure-proof costs are
unknown, not zero.

Put $z(t)=U(0,t)$ and let $w$ be the antidiagonal permutation matrix.
The standard coordinates are the matrices in Bray--B\"a\"arnhielm,
\emph{A new method for recognising Suzuki groups}, arXiv:1706.05697v1,
equations (1)--(5), physical page 1; the exact local primary is bound in
the dependency record. They are
\[
z(t)=\begin{pmatrix}1&0&0&0\\0&1&0&0\\t&0&1&0\\t^\theta&t&0&1\end{pmatrix},
\qquad
h(\lambda)=\operatorname{diag}(\lambda^{\theta+1},\lambda,
\lambda^{-1},\lambda^{-\theta-1}).
\]
The root group centralizes all $z(t)$. Conjugation by $h(\lambda)$ scales
$t$ by $\lambda^{-\theta-2}$. This exponent is coprime to $2^n-1$:
write $n=2e+1$, discard the factor two, and reduce $2^{2e+1}-1$ modulo
$2^e+1$ to obtain one. Thus scaling runs through the whole multiplicative
field group. The scaling parameter belongs to $\mathbb F_{2^d}$ precisely
when $\lambda$ does, since a power automorphism of a cyclic group preserves
each of its subgroups. The twist restricts to the subfield twist because
$(n+1)/2\equiv(d+1)/2\pmod d$ for the odd extension degree $n/d$.
We use ordinary Sylow conjugacy for finite groups.

\section{Uniform seeds without the source conjugation manipulation}
\begin{lemma}[Trace and local seeds]\label{seed}
The $q-1$ nonzero traces distinguish all nonidentity odd-order conjugacy
classes in $G$. Suppose $n=3a$, $Q=2^a$, and $\tau=F^a$ or $F^{2a}$.
For each ovoid point $P$ moved by $\tau$, the products
$\rho\tau(\rho)$, as $\rho$ ranges over the $q-1$ involutions fixing $P$,
represent each such class exactly once.
If $h>1$ is an odd element order of $G$ and $c_h$ its number of conjugacy
classes, the number of involutions with $|\rho\tau(\rho)|=h$ is
\[
 c_h(q^2-Q^2)=c_h\,|\Sz(Q)|(Q+1).
\]
In particular $c_5=1$.
\end{lemma}
\begin{proof}
The matrices $wz(t)$ have trace $t^\theta$, hence distinct nonzero traces.
They are nonidentity and of odd order by the retained even-order trace
property. The number of nonidentity odd conjugacy classes, obtained from
the three tori and their freely acting normalizer quotients, is
\[
 \frac{q-2}{2}+\frac{q+\sqrt{2q}}4+
 \frac{q-\sqrt{2q}}4=q-1.
\]
Consequently the exhibited traces exhaust and distinguish these classes.

Conjugate $P,\tau(P)$ to $\infty,0$. In these coordinates $\tau$ is a
semilinear map with field component $t\mapsto t^{Q^\epsilon}$,
$\epsilon=1$ or $2$, taking $\infty$ to $0$. After composing with $w$,
its inner part fixes $\infty$, so belongs to $U H$. Its action on the
central root involutions therefore has the form
\[
 \tau(z(t))=wz(c t^{Q^\epsilon})w\qquad(c\in\mathbb F_q^*).
\]
Direct matrix multiplication gives
\[
 \operatorname{tr}(z(b)wz(c')w)=(bc')^\theta,
 \quad\text{hence}\quad
 \operatorname{tr}(\rho\tau(\rho))=(c t^{1+Q^\epsilon})^\theta.
\]
Both $1+Q$ and $1+Q^2$ are coprime to $Q^3-1$. For the first gcd,
$Q\equiv-1$ gives $Q^3-1\equiv-2$. For the second,
$Q^2\equiv-1$ and $Q^3\equiv1$ give $Q\equiv-1$ and then
$2\equiv0$. Since $Q$ is even, both gcds are odd and therefore one.
Thus these traces traverse all nonzero field elements. This proves the local assertion and, in particular,
repairs source Lemma 3.1 directly without its conjugation calculation.
At a fixed ovoid point $\rho$ and $\tau(\rho)$ lie in the same elementary
abelian center, so cannot have product of odd order $h>1$.
There are $Q^2+1$ fixed ovoid points, giving the count. The identity
$q^2-Q^2=|\Sz(Q)|(Q+1)$ follows from $q=Q^3$ and
$|\Sz(Q)|=Q^2(Q^2+1)(Q-1)$, also for $Q=2$.
Finally $5$ divides a nonsplit torus; its four nonidentity elements fuse
under the order-four normalizer quotient, so $c_5=1$.
\end{proof}

For $n=3a$, every automorphism of order three of $G$ is conjugate by $G$ to $F^a$ or
$F^{2a}$. Indeed $\Aut(G)=G\rtimes\gen{F}$, and the outer image of such an
automorphism is nontrivial because $3\nmid|G|$. It lies in the order-three
subgroup of the cyclic outer group. In its inverse image
$G\rtimes C_3$, the order-three subgroups are Sylow subgroups, so are
conjugate. The conjugator can be chosen in $G$ by removing its field
component; the quotient image selects $F^a$ or $F^{2a}$. Consequently
the local seed statement also covers the source's arbitrary order-three
automorphism and every point moved by it. The count above is invariant
under that conjugation. The identity order $h=1$ is excluded: at a moved
point the two involutions are distinct.

\section{Subfield descent and conjugacy}
\begin{lemma}[Descent of a cyclic marking]\label{descent}
Let $\rho_i=\tau^i(\rho_0)$ be involutions whose pair products all have
odd order $h>1$. If their three cyclic rotation subgroups are distinct,
then $H=\langle\rho_0,\rho_1,\rho_2\rangle$ is a conjugate of some
$S_d$, $d\mid n$, $d>1$. The same conclusion holds automatically when
$h\equiv2\pmod3$.
\end{lemma}
\begin{proof}
The involutions have distinct fixed points, since commuting root-center
involutions have product of order at most two. A Borel cannot contain
$H$. Within a torus normalizer, all odd elements lie in its unique cyclic
torus, hence all three order-$h$ rotation subgroups would coincide.
This excludes that maximal-subgroup type under the first hypothesis.

For the second hypothesis, torus-normalizer containment would put the
three involutions in the same dihedral group $D_{2h}$: from
$\rho_0\rho_1,\rho_1\rho_2\in C_h$ follows $\rho_2\in\rho_0C_h$.
Write $x=\rho_0\rho_1$ and $\tau(x)=x^u$. Then
$\tau^2(x)=x^{u^2}$ and
\[
 1=(\rho_0\rho_1)(\rho_1\rho_2)(\rho_2\rho_0)
   =x^{1+u+u^2}.
\]
An odd $h\equiv2\pmod3$ has a prime divisor $\ell\equiv2\pmod3$.
But the cube map is bijective modulo $\ell$, so
$u^2+u+1=0$ would force $u=1$ and then $3=0$, impossible.
Thus the torus-normalizer case is excluded here too.
If $H$ is proper, put it in a maximal subfield subgroup and repeat the
same argument there. Each field degree strictly decreases. Eventually
$H$ equals a Suzuki subfield group. The terminal $S_1=C_5:C_4$ is itself
a torus normalizer, so cannot occur. Nested subfield embeddings show that
the final subgroup is a conjugate of a canonical $S_d$ in $G$.
\end{proof}

\begin{lemma}[All normalized subfields counted]\label{subfields}
The normalizer of every canonical $S_d$, including $d=1$, in $G$ is $S_d$.
For $\tau=F^a$, the $\tau$-stable conjugates of $S_d$ form one orbit under
$C_G(\tau)=S_a$, and their number is
\[
 \frac{|S_a|}{|S_{\gcd(d,a)}|}.
\]
If a cyclic marking as in Lemma~\ref{descent} generates a proper subfield,
then $3\mid d$ and $\gcd(d,a)=d/3$. In particular $n/d\ge5$.
\end{lemma}
\begin{proof}
The set of ambient fixed points of the involutions of $S_d$ is exactly
its subovoid. A normalizing element preserves this set. Use the doubly
transitive action of $S_d$ on that subovoid to adjust the element to fix
$\infty,0$. It is now in the ambient split torus. Its scaling preserves
$z(\mathbb F_{2^d})$, so the scaling parameter, hence the torus element,
lies in the subfield. Undoing the adjustment gives $N_G(S_d)=S_d$.
This also works for the sharply doubly transitive group $S_1$ on five points.

Work in $G\rtimes\langle\tau\rangle$. If $S_d^g$ is $\tau$-stable,
the conjugate of $\tau$ by $g^{-1}$ normalizes $S_d$ and lies in
$S_d\rtimes\langle\tau\rangle$, by the normalizer just proved.
Since $3$ does not divide $|S_d|$, its cyclic subgroup of order three is
conjugate to $\langle\tau\rangle$ by Sylow conjugacy. The conjugator may
be chosen in $S_d$, by removing its $\tau$ component. Conjugation preserves
the image in the quotient $C_3$, so it takes the element to $\tau$, not
$\tau^2$. It follows that $S_d^g$ is conjugate to $S_d$ by an element of
$C_G(\tau)$. Its stabilizer there is precisely
$S_d\cap C_G(\tau)=S_{\gcd(d,a)}$, proving the formula.

If $d\mid a$, every such stable conjugate is fixed pointwise by $\tau$,
which is incompatible with the marking. For $d\mid n=3a$ the only other
possibility is $3\mid d$, $\gcd(d,a)=d/3$. If $d<n$, the odd integer
$n/d$ is greater than one and not divisible by three, hence at least five.
\end{proof}

\section{A repaired all-parameter generation statement}
\begin{theorem}[Uniform existence, not universal generation]\label{uniform}
Let $n=3a>1$ be odd, $G=\Sz(2^n)$, and $\tau=F^a$.
For every odd element order $h>1$ with $h\equiv2\pmod3$, there is an
involution $\rho$ such that
\[
 |\rho\tau(\rho)|=h,\qquad
 \langle\rho,\tau(\rho),\tau^2(\rho)\rangle=G.
\]
In particular this holds for $h=5$ for every such $n$.
The number of fully generating seeds is at least
\[
 c_h |\Sz(2^a)|\left(2^a+1-
 \sum_{u=1}^{\lfloor a/5\rfloor}(2^u+1)\right)>0.
\]
The same assertions hold for the inverse field automorphism.
\end{theorem}
\begin{proof}
By Lemma~\ref{seed} there are $c_h|S_a|(2^a+1)$ seeds. By
Lemma~\ref{descent}, a nongenerating seed generates a proper Suzuki
subfield. A stable subfield of degree $d=3u$ can contribute only if
$u\le a/5$, by Lemma~\ref{subfields}. Write $s=2^u$; the restriction
of $\tau$ to $S_d$ is one of its two order-three field automorphisms.
If $h$ occurs in $S_d$, its number of conjugacy classes there equals the
ambient $c_h$. Indeed it is $\varphi(h)/2$ in a split torus and
$\varphi(h)/4$ in a nonsplit torus. Split type is preserved under a field
extension. For nonsplit type, $r^2\equiv-1\pmod h$, $r=2^d$; every
extension degree is odd, so $q^2\equiv-1\pmod h$ and $h>1$ cannot divide
$q-1$. Thus its class count is unchanged. If $h$ does not occur, this
subfield contributes no seed.

Each stable conjugate contributes at most $c_h(s^6-s^2)$ seeds. Multiply
by the exact conjugate count $|S_a|/|\Sz(s)|$ and use
\[
 \frac{s^6-s^2}{s^2(s^2+1)(s-1)}=s+1.
\]
Sum over proper active degrees, bounded by all integers
$1\le u\le b=\lfloor a/5\rfloor$. This is an upper bound on a union;
intersections only decrease it. Subtracting gives the displayed result.
For $b=0$ there is no subtraction. For $b\ge1$,
\[
 \sum_{u=1}^b(2^u+1)=2^{b+1}-2+b
 <2^{b+2}\le2^{5b}\le2^a,
\]
since $b\le2^b$ and $b+2\le5b$. This proves strict positivity.
\end{proof}

\section{The literal counterexample remains a counterexample}
The preceding theorem repairs existence; it does not repair the printed
unrestricted implication. Apply Theorem~\ref{uniform} inside $K=\Sz(8)$
with its field automorphism $F^2$, obtaining an involution $\rho$ whose
cyclic triple generates $K$ and whose pair products have order five.
This is a written existence proof, not a new finite search premise.
Embed the standard matrices in $G=\Sz(2^{15})$ and use $\tau=F^5$.
The twist $2^8$ restricts on $\mathbb F_8$ to $2^2$, the degree-three
Suzuki twist, so this is the standard subfield inclusion. The restriction
of $\tau$ is $F^2$; hence the same three involutions still generate
precisely $K$, not $G$. Their cyclic order-five subgroups are distinct:
if all coincided, the generated group would be contained in its torus
normalizer, contrary to generation of $K$. Cyclic symmetry means equality
of any pair implies equality of all three. The first involution's unique
ovoid point is moved by $\tau$, since otherwise its product with its
image would have order at most two. Thus every printed hypothesis in the
second assertion of Theorem 3.2 holds, while full generation fails.
The source's maximal-stable-subfield inference overlooks the index-five
subfield here. This proof also supplies the exact abstract counterexample
independently of the earlier, accepted order-seven L122 matrix marking.

The three supports have five points inside the 65-point subovoid. They
cannot all coincide, since then $K$ would stabilize a proper nonempty
subset of its transitive subovoid. Hence the same example refutes the
unrestricted ambient-group conclusion of Corollary 3.5 and the ``any
triangle'' assertion of Corollary 3.7, even for order five. The application
will prove its triangle is proper as well. Its chamber graph on the
ambient group has components the cosets of $K$, so the unconditional
ambient residual-connectivity assertion also requires correction.
None of this refutes the existential main Theorem 3.16.

\paragraph{Application and scope.}
The accompanying application proves prime-order properness, the rank-three product identity,
the exact source chamber-system identification, and the full correlation
and no-duality conclusions. Its common coset-recognition dependency is
supplied in full. The source's universal auxiliary statements receive
separate corrected or refuted dispositions. This unreviewed draft assigns
no complete-source credit.
\end{document}
